\documentclass[tikz,border=6pt]{standalone}
\usepackage[fontset=fandol]{ctex}
\usepackage{tikz}
\usetikzlibrary{arrows.meta, positioning, fit, calc}

\begin{document}
\begin{tikzpicture}[
  font=\small,
  box/.style={draw, rounded corners, align=center, minimum width=3.4cm, minimum height=1.05cm, fill=blue!5},
  obs/.style={draw, rounded corners, align=center, minimum width=3.2cm, minimum height=1.05cm, fill=green!6},
  action/.style={draw, rounded corners, align=center, minimum width=3.4cm, minimum height=1.05cm, fill=orange!8},
  arrow/.style={-Latex, thick},
  dashedarrow/.style={-Latex, thick, dashed}
]

\node[box] (h) {竞争性假设\\\footnotesize（回路结构、功能解释）};
\node[box, right=1.8cm of h] (m) {模型推演\\\footnotesize（输出、时序、分歧点）};
\node[action, right=1.8cm of m] (e) {关键实验设计\\\footnotesize（测哪里、刺激哪里）};

\node[obs, below=1.5cm of m] (d) {观测数据\\\footnotesize（多通道、连续、长时间窗）};
\node[obs, below=1.5cm of e] (p) {in silico  扰动\\\footnotesize（刺激、切换、比较）};
\node[obs, below=1.5cm of h] (s) {多尺度同步验证\\\footnotesize（底层到上层的转移函数）};

\draw[arrow] (h) -- (m);
\draw[arrow] (m) -- (e);
\draw[arrow] (e) |- (d);
\draw[arrow] (d) -| (h);
\draw[arrow] (m) -- (p);
\draw[arrow] (p) -- (e);
\draw[dashedarrow] (s) -- (d);
\draw[dashedarrow] (s) -- (h);

\node[draw, rounded corners, fit=(h)(m)(e)(d)(p)(s), inner sep=0.45cm, line width=0.8pt, color=gray!70] {};

\end{tikzpicture}
\end{document}
